\appendix
\part{Anhang}

\newcommand{\lit}[3]{\vspace{5pt}\noindent\textbf{[#1]} #2\par\nopagebreak\vspace{1pt}{\small\itshape\color{grau} #3}\par\vspace{2pt}}

\chapter{Formelsammlung und Übersicht}

\section{Begriffe und Formeln}

\begin{merksatz}
\boxtitel{Teil I bis III}
\begin{center}
{\small
\renewcommand{\arraystretch}{1.6}
\begin{tabular}{@{}P{4.6cm}P{4.0cm}@{}}
\toprule
$\hat\theta = \frac1n\sum h(X_i)$ & Monte-Carlo-Schätzer (1.1)\\
$\operatorname{SE} = \sigma/\sqrt n$ & Wurzelregel (1.2)\\
$B_c = aB_{c-1}/(c+aB_{c-1})$ & Erlang-Rekursion (1)\\
$x_{k+1} = (ax_k+c)\bmod m$ & LCG (2.1)\\
$\le (d!\,m)^{1/d}$ Ebenen & Marsaglia (2.3)\\
$X = F^{-1}(U)$ & Inversion (2.4)\\
$\sqrt{-2\ln U_1}\cos 2\pi U_2$ & Box--Muller (2.5)\\
$n^{-2/d}$ gegen $n^{-1/2}$ & Fluch der Dimension (3.1)\\
$\tau = 1+2\sum\rho_k$, $n/\tau$ & Autokorrelationszeit (4.1)\\
$s_{\text{Block}}/\sqrt b$ & Blockmittel (4.2)\\
$\operatorname{Var}[A-B] = \ldots - 2\operatorname{Cov}$ & gemeinsame Zahlen (4.3)\\
$\tfrac12(g(U)+g(1-U))$, Faktor $1+\rho$ & antithetisch (5.1)\\
$Y-\beta^*(C-c)$, $1-\rho^2$ & Kontrollvariable (5.2)\\
$\operatorname{Var}[\mathbb E[Y|Z]]\le\operatorname{Var}[Y]$ & Konditionierung (5.3)\\
$w = f/g$, $g^*\propto|h|f$ & Wichtigkeit (6.1), (6.2)\\
$e^{\mu^2}P(Z>c+\mu)$ & zweites Moment, Verschiebung\\
\bottomrule
\end{tabular}}
\end{center}
\begin{center}
{\small
\renewcommand{\arraystretch}{1.6}
\begin{tabular}{@{}P{4.6cm}P{4.0cm}@{}}
\toprule
$(\sum w)^2/\sum w^2$ & wirksames $n$\\
$|\bar h - \int h| \le V(h)D^*_n$ & Koksma--Hlawka (6.3)\\
$\min\{1, \tfrac{\pi(x')q(x|x')}{\pi(x)q(x'|x)}\}$ & Metropolis--Hastings (7.1)\\
$\pi(x)P(x\to x') = \pi(x')P(x'\to x)$ & detaillierte Bilanz (7.2)\\
$2{,}38/\sqrt d$, $23\,\%$ & optimale Schrittweite\\
$\hat R = \sqrt{\hat V/W}$ & Gelman--Rubin (8.1)\\
$H = U + \tfrac12 p^\top p$, Sprungverfahren & Hamilton (9.1), (9.2)\\
$\Delta t^{1/2}$, $\Delta t$, $\Delta t$ & Euler, Milstein, schwach (9.4)\\
$N_\ell\propto\sqrt{V_\ell/C_\ell}$ & Mehrebenen (9.5)\\
\bottomrule
\end{tabular}}
\end{center}
\end{merksatz}

\begin{merksatz}
\boxtitel{Teil IV und V}
\begin{center}
{\small
\renewcommand{\arraystretch}{1.6}
\begin{tabular}{@{}P{4.6cm}P{4.0cm}@{}}
\toprule
(Zeit, Rang, Art) & Ereignisliste (10.1)\\
$L = \lambda_{\text{eff}}\,\bar W$ & Little (10)\\
$\tfrac12\Delta t$ zu lang & Zeitschritt-Verzerrung (10)\\
$\rho^{\sqrt{2(c+1)}-1}/(c(1-\rho))\cdot\tfrac{1+c_s^2}{2}\,S$ & Sakasegawa (11.1)\\
$(1-r)/(1-r^{\ell})$, $r = (1-p)/p$ & Spielerruin (11)\\
$n_m/(N_0R^{m-1})$ & Aufspaltung (11.2)\\
Vorhersage, Gewichtung, Neuziehung & Teilchenfilter (12.1)\\
$K = (P+q)/(P+q+r)$ & Kalman-Gewinn (12)\\
Erhaltung, Grenzen, Formel, Unempfindlichkeit, zweites Programm & Prüfleiter (13.1)\\
$\int\mathbf 1\{g(\theta)\le z\}\,p(\theta|\text{D})\,d\theta$ & A-posteriori-Vorhersage (14.1)\\
$S_j = \operatorname{Var}[\mathbb E[g|\theta_j]]/\operatorname{Var}[g]$ & Sobol-Index (14.2)\\
$(a/B)\,dB/da$ & Elastizität (14)\\
\bottomrule
\end{tabular}}
\end{center}
\end{merksatz}

\section{Die Landkarte}

\begin{merksatz}
\boxtitel{Welcher Satz worauf beruht}
\begin{center}
{\small
\renewcommand{\arraystretch}{1.5}
\begin{tabular}{@{}P{3.4cm}P{5.2cm}@{}}
\toprule
\textbf{Satz} & \textbf{beruht auf}\\
\midrule
1.2 Wurzelregel & Gesetz der großen Zahlen, zentraler Grenzwertsatz (Band 17)\\
2.3 Marsaglia & Gitter in $\mathbb Z^d$ (Band 1, 7)\\
2.4 Inversion & Verteilungsfunktion (Band 15)\\
2.6 Verwerfung & bedingte Wahrscheinlichkeit (Band 15)\\
3.1 Fluch & Fehler der Trapezregel (Band 11)\\
4.1 Autokorrelationszeit & Satz 10.4 (Band 26)\\
5.1 antithetisch & Chebyshev-Ungleichung für Funktionen (Band 15)\\
5.2 Kontrolle & Regression (Band 26)\\
5.3 Konditionierung & Turmregel, Varianzzerlegung (Band 25)\\
6.2 Wichtigkeit & Dichtewechsel (Band 15)\\
6.3 Koksma--Hlawka & partielle Integration; mehrdimensional Auslassung 2\\
7.2 Metropolis & Sätze 3.5 und 3.9 (Band 25)\\
8.2 $\hat R$ & Varianzzerlegung (Band 17)\\
9.2 Hamilton & Liouville (Band 24, Auslassung 3); 7.2\\
9.4 Ordnung & Satz 8.6 (Band 25); Milstein Auslassung 4\\
9.5 Mehrebenen & Lagrange (Band 12); 5.2\\
\bottomrule
\end{tabular}}
\end{center}
\begin{center}
{\small
\renewcommand{\arraystretch}{1.5}
\begin{tabular}{@{}P{3.4cm}P{5.2cm}@{}}
\toprule
\textbf{Satz} & \textbf{beruht auf}\\
\midrule
11.1 Sakasegawa & Erlang C, Pollaczek--Chintschin (Band 25); Näherung\\
12.2 Teilchenfilter & 6.1; Satz von Bayes (Band 21)\\
13.1 Prüfleiter & kein Satz --- eine Praxis\\
14.3 Sobol & orthogonale Zerlegung (Band 20, Auslassung 5)\\
\bottomrule
\end{tabular}}
\end{center}
\end{merksatz}

\section{Was nicht bewiesen wurde}

\begin{merksatz}
\boxtitel{Fünf Auslassungen}
\begin{center}
{\small
\renewcommand{\arraystretch}{1.5}
\begin{tabular}{@{}P{3.8cm}P{1.0cm}P{3.2cm}@{}}
\toprule
\textbf{Ergebnis} & \textbf{Kap.} & \textbf{Wo es steht}\\
\midrule
Volle Periode des LCG (Hull--Dobell) & $2$ & Band 1, Kap.~9; Knuth\zit{12}\\
Koksma--Hlawka, mehrdimensional & $6$ & Band 20; Niederreiter\zit{21}\\
Satz von Liouville & $9$ & Band 24; Neal\zit{20}\\
Starke Ordnung des Milstein-Verfahrens & $9$ & Band 25; Kloeden--Platen\zit{11}\\
Sobol-Zerlegung & $14$ & Band 20; Saltelli\zit{26}\\
\bottomrule
\end{tabular}}
\end{center}

\begin{center}
\begin{minipage}{0.88\textwidth}
\kursiv{\textbf{Alle übrigen Sätze des Bandes sind vollständig bewiesen} --- darunter die Wurzelregel, der Ebenensatz für RANDU, Inversion und Verwerfung, der Fluch der Dimension, die Varianzformeln der antithetischen und der Kontrollvariablen, die optimale Vorschlagsdichte, die detaillierte Bilanz des Metropolis-Algorithmus, die Annahme $1$ des Gibbs-Samplers, die Volumenerhaltung des Sprungverfahrens, die Teleskopzerlegung mit optimaler Pfadverteilung und die Erwartungstreue des Teilchenfilters. Die Näherung von Sakasegawa ist eine Interpolationsformel ohne Beweis und wurde als solche gekennzeichnet.}
\end{minipage}
\end{center}
\end{merksatz}

\section{Zeittafel}

\begin{merksatz}
\boxtitel{Von Buffon bis zur Kopplung}
\begin{center}
{\small
\renewcommand{\arraystretch}{1.42}
\begin{tabular}{@{}P{1.0cm}P{3.4cm}P{4.2cm}@{}}
\toprule
$1777$ & Buffon & die Nadel: $\pi$ durch Zufall\\
$1909$ & Erlang & Verlustformel für Telefonzentralen\\
$1916$ & Weyl & Gleichverteilung modulo eins\\
$1946$ & Ulam, von Neumann & Monte Carlo in Los Alamos\\
$1949$ & Lehmer & linearer Kongruenzgenerator\\
$1951$ & Kahn, Harris & Wichtigkeitsstichprobe und Aufspaltung\\
$1953$ & Metropolis u.\,a. & der Algorithmus\\
$1955$ & Maruyama & Euler für stochastische Gleichungen\\
$1956$ & Hammersley, Morton & antithetische Variablen\\
$1958$ & Box, Muller & normalverteilte Zahlen\\
$1960$ & Kalman & der Filter\\
$1961$ & Little; Gordon; Dahl, Nygaard & $L = \lambda W$; GPSS; Simula\\
$1967$ & Sobol & Folgen niedriger Diskrepanz\\
$1968$ & Marsaglia & \enquote{Random numbers fall mainly in the planes}\\
$1970$ & Hastings & Verallgemeinerung auf unsymmetrische Vorschläge\\
$1974$ & Milstein & Verfahren erster Ordnung\\
\bottomrule
\end{tabular}}
\end{center}
\begin{center}
{\small
\renewcommand{\arraystretch}{1.42}
\begin{tabular}{@{}P{1.0cm}P{3.4cm}P{4.2cm}@{}}
\toprule
$1984$ & Geman, Geman & Gibbs-Sampler\\
$1987$ & Duane u.\,a. & Hybrid Monte Carlo\\
$1990$ & Gelfand, Smith; Sobol & Bayes wird rechenbar; Varianzzerlegung\\
$1992$ & Gelman, Rubin & $\hat R$\\
$1993$ & Gordon, Salmond, Smith & Bootstrap-Filter\\
$1995$ & Owen & verwürfelte Netze\\
$1997$ & Roberts, Gelman, Gilks & $0{,}234$\\
$1998$ & Matsumoto, Nishimura & Mersenne-Twister\\
$2008$ & Giles & Mehrebenen-Monte-Carlo\\
$2010$ & Andrieu, Doucet, Holenstein & Teilchen-MCMC\\
$2011$ & Salmon u.\,a. & Zählergeneratoren (Philox)\\
$2014$ & Hoffman, Gelman; O'Neill & NUTS; PCG\\
$2020$ & Jacob, O'Leary, Atchadé & unverzerrte Ketten durch Kopplung\\
\bottomrule
\end{tabular}}
\end{center}
\end{merksatz}

\section{Glossar}

\begin{merksatz}
\boxtitel{Begriffe von A bis Z}
\begin{center}
{\small
\renewcommand{\arraystretch}{1.3}
\begin{longtable}{@{}P{3.6cm}P{5.0cm}@{}}
\toprule
\textbf{Begriff} & \textbf{Bedeutung}\\
\midrule
\endhead
Aleatorisch & Zufall, der bleibt (14)\\
Antithetisch & $U$ und $1-U$ als Paar (5.1)\\
Aufspaltung & gute Pfade kopieren (11.2)\\
Autokorrelationszeit & $1+2\sum\rho_k$ (4.1)\\
Blockmittel & Standardfehler ohne $\rho_k$ (4.2)\\
Box--Muller & zwei Gleichverteilte, zwei Normale (2.5)\\
Degeneration & alle Gewichte auf einem Teilchen (12)\\
Detaillierte Bilanz & Fluss hin $=$ Fluss zurück (7.2)\\
Diskrepanz & wie ungleichmäßig eine Punktmenge ist (6.3)\\
Einschwingphase & Anfang verwerfen (4)\\
Epistemisch & Zufall, der schrumpft (14)\\
Ereignisliste & nach Zeit geordnete Halde (10.1)\\
Erlang B & Verlustformel, unempfindlich (1)\\
Fluch der Dimension & Gitter brauchen $n^{d}$ (3.1)\\
Gemeinsame Zufallszahlen & gleicher Zufall für Varianten (4.3)\\
Gibbs-Sampler & bedingt ziehen, nie ablehnen (7.3)\\
Hamilton-Monte-Carlo & Impuls und Bahn statt Zufallsweg (9.1)\\
Inversion & $F^{-1}(U)$ (2.4)\\
Kontrollvariable & bekannten Erwartungswert abziehen (5.2)\\
Konditionierung & $\mathbb E[Y|Z]$ statt $Y$ (5.3)\\
LCG & $ax+c \bmod m$ (2.1)\\
Mehrebenen & Teleskopsumme über Gitter (9.5)\\
Metropolis--Hastings & vorschlagen, annehmen mit $\alpha$ (7.1)\\
Milstein & Euler plus $\tfrac12 bb'(\Delta W^2-\Delta t)$ (9.3)\\
Neuziehung & Teilchen nach Gewicht vervielfachen (12.1)\\
$\hat R$ & zwischen/innerhalb der Ketten (8.1)\\
RANDU & fünfzehn Ebenen (2.3)\\
Sobol-Folge & Punkte niedriger Diskrepanz (6)\\
Sobol-Index & Varianzanteil einer Eingabe (14.2)\\
Sprungverfahren & symplektischer Integrator (9.2)\\
Startwert & Anfang der Zahlenfolge (2, 13)\\
Teilchenfilter & Wichtigkeit Tag für Tag (12.1)\\
Verwerfung & Vorschlag unter die Dichte (2.6)\\
Wichtigkeitsstichprobe & aus $g$ ziehen, mit $f/g$ wichten (6.1)\\
Wirksames $n$ & $n/\tau$ oder $(\sum w)^2/\sum w^2$ (4, 6)\\
Wurzelregel & Fehler $\sigma/\sqrt n$ (1.2)\\
Zählergenerator & Zahl aus Zählerstand (2, 13)\\
\bottomrule
\end{longtable}}
\end{center}
\end{merksatz}
